Here we will intruduce the differnet types of relevant faults and assumptions.
\todo[inline]{I think we could do some cool graphics here, describing the faults.}
\begin{itemize}
    \item \textbf{Actuator faults:} Loss-of-effectiveness (LoE), stuck-on/off, bias, saturation.

    \item \textbf{Sensor faults:} bias, drift, dropout, complete outage.

    \item \textbf{Contraints:} actuator limits (torque, thrust, rate)

    \item \textbf{Inertia and velocity sensing assumption:} Most spacecraft FTC studies consider the inertia tensor uncertain, since on-orbit conditions (propellant slosh, deployable structures) can shift the mass distribution. In contrast, the present air-bearing testbed has a rigid body with fixed, measurable inertia, so inertia is treated as known.

Another spacecraft-specific theme is velocity sensing. Many CubeSats exclude gyros to save cost and power, relying only on attitude sensors (e.g., star trackers, magnetometers). This motivates “velocity-free FTC” designs, where angular rates are reconstructed with observers. The testbed provides direct rate measurements, but the velocity-free paradigm remains an important thread in spacecraft FTC literature.
\end{itemize}


